% ภาคผนวก C: อภิธานศัพท์ (คำสั่ง \appendix อยู่ที่ appendix-a แล้ว)
% ============================================================================
\mychapter{อภิธานศัพท์}

\noindent อภิธานศัพท์นี้รวบรวมคำศัพท์ที่ใช้ในคู่มือและในระบบ

\begin{small}
\renewcommand{\arraystretch}{1.25}
\begin{longtable}{L{4.5cm}L{10cm}}
\toprule
\textbf{คำศัพท์} & \textbf{คำอธิบาย} \\
\midrule
จัดอัตโนมัติ & บทบาทของไฟล์ที่ระบบจะจัดเวลาสอบให้โดยอัตโนมัติ \\
จัดไปแล้ว & บทบาทของไฟล์ที่ระบบจะคงเวลาสอบเดิมจากไฟล์ไว้ \\
ตอนเรียน & กลุ่มย่อยของวิชาที่เปิดสอนแยกกัน \\
กลุ่มนักศึกษา & กลุ่มของนักศึกษาที่ต้องเข้าสอบร่วมกัน \\
รหัสวิชา & ตัวระบุวิชา เก้าหลัก \\
ล็อกเวลา & สถานะของเวลาสอบที่ผู้ใช้กำหนดเอง \\
ต้นทางกำหนด & จำนวนวิชาที่มีเวลาสอบมาจากไฟล์ต้นทาง (แสดงบนหน้าภาพรวม) \\
ล็อกเอง & จำนวนวิชาที่ผู้ใช้ล็อกเวลาเอง (แสดงบนหน้าภาพรวม) \\
สอบเต็มวัน & วิชาที่ต้องสอบเต็มวันตามช่วงที่ตั้งค่าไว้ \\
สอบตรงกัน & วิชาที่จัดสอบพร้อมกันกับวิชาอื่น \\
ข้อมูลต้นทาง & ไฟล์ข้อมูลที่นำเข้าสู่ระบบ \\
ตารางสอบ & ผลลัพธ์ของการจัดเวลาสอบ \\
ปฏิทินสอบ & การแสดงตารางสอบในรูปแบบปฏิทิน \\
ข้อผิดพลาด & ปัญหาระดับรุนแรงที่ขัดขวางการจัดตาราง \\
คำเตือน & ปัญหาระดับเบาที่ควรทราบ \\
ฉบับ & สถานะของโครงการที่บันทึกไว้บนคลาวด์ ณ เวลาหนึ่ง การบันทึกแต่ละครั้งสร้างฉบับใหม่ \\
ลิงก์แชร์ & ลิงก์สำหรับเปิดโครงการที่แชร์ \\
สิทธิ์ & ระดับการเข้าถึงของลิงก์แชร์ \\
สร้างอัตโนมัติ & ค่าคอลัมน์ ที่มา ในไฟล์ส่งออก เมื่อระบบจัดเวลาเอง (เรียกว่า จัดอัตโนมัติ ในหน้าจอ) \\
กำหนดเอง & ค่าคอลัมน์ ที่มา ในไฟล์ส่งออก เมื่อผู้ใช้แก้ไขเวลาเอง (สถานะ ล็อกเวลา) \\
นำเข้า & ค่าคอลัมน์ ที่มา ในไฟล์ส่งออก เมื่อเวลามาจากไฟล์ต้นทาง (เรียกว่า ต้นทางกำหนด) \\
\bottomrule
\end{longtable}
\end{small}
